% Generated draft: requires mathematical and source review.
\documentclass[11pt]{article}
\usepackage[margin=25mm]{geometry}
\usepackage{fontspec}
\setmainfont{Latin Modern Roman}
\usepackage{amsmath,amssymb,mathtools,unicode-math}
\setmathfont{Latin Modern Math}
\usepackage{enumitem,needspace,xurl,hyperref}
\setlength{\parindent}{0pt}
\setlength{\parskip}{5pt}
\setlength{\emergencystretch}{3em}
\newcommand{\N}{\mathbb{N}}
\newcommand{\Z}{\mathbb{Z}}
\newcommand{\Q}{\mathbb{Q}}
\newcommand{\R}{\mathbb{R}}
\newcommand{\C}{\mathbb{C}}
\newcommand{\F}{\mathbb{F}}
\newcommand{\A}{\mathbb{A}}
\newcommand{\PP}{\mathbb{P}}
\newcommand{\E}{\mathbb{E}}
\newcommand{\eps}{\varepsilon}
\newcommand{\dd}{\,\mathrm{d}}
\begin{document}
\section*{UnsolvedMath ID: 30003850}
\textbf{Source problem code:} catalogue entry
\subsection{Definitions and mathematical statement}
The supplied record merges two independent numbered questions and marks the extracted statement as unrecoverable. The following separates its mathematical content while retaining unresolved qualifications; it is not a verified reconstruction of either original numbered question.

A discrete group \(G\) is amenable if there exists a finitely additive function \(m:\mathcal P(G)\to[0,1]\) with \(m(G)=1\) and \(m(gA)=m(A)\) for every \(g\in G\) and every subset \(A\subseteq G\). It is nonamenable otherwise. A \(K(G,1)\) is a connected CW complex with fundamental group isomorphic to \(G\) and contractible universal cover. A finite \(K(G,1)\) has finitely many cells in total. A group has finite virtual cohomological dimension, in the usual virtually torsion-free sense used here, if it has a torsion-free subgroup \(H\leq G\) of finite index with \(\operatorname{cd}_{\Z}(H)<\infty\), where \(\operatorname{cd}_{\Z}(H)\) is the projective dimension of the trivial \(\Z H\)-module \(\Z\).

Let \(F_2\) denote the free group on two generators. Conditional on interpreting the source phrase ``no free subgroups'' as ``no nonabelian free subgroups,'' the first question asks whether
\[
\exists G\;\bigl(G\text{ is nonamenable},\quad G\text{ admits a finite }K(G,1),\quad F_2\not\hookrightarrow G\bigr),
\]
or, as a weaker target, whether
\[
\exists G\;\bigl(G\text{ is nonamenable},\quad G\text{ has finite virtual cohomological dimension},\quad F_2\not\hookrightarrow G\bigr).
\]
Here \(F_2\not\hookrightarrow G\) means that there is no injective group homomorphism from \(F_2\) to \(G\); equivalently, \(G\) contains no nonabelian free subgroup. The original wording does not specify the rank of the excluded free subgroups. Excluding all nontrivial free subgroups would also exclude infinite cyclic subgroups and would be a different condition. The displayed formulations are therefore explicitly conditional interpretations.

The second extracted question asks whether Wajnryb--Witowicz-regular relations for the group denoted \(B\), described in the source as the amenable Brin--Navas group, can be used to establish such relations for Thompson's group \(F\). Here \(F\) is the group, under composition, of increasing piecewise-linear homeomorphisms of \([0,1]\) with finitely many breakpoints, all breakpoints dyadic rational, and all slopes integral powers of \(2\). Such maps fix both endpoints. The supplied record does not define \(B\), Wajnryb--Witowicz regularity, the relevant relations, or their quantifiers. Consequently this second question cannot be converted into a complete mathematical predicate from the supplied evidence, and no connection between it and the preceding existence questions is asserted.

For interpreting the supplied literature assessment, type \(F_\infty\) means that a group admits a \(K(G,1)\) with finite \(n\)-skeleton for every nonnegative integer \(n\). This differs from having a finite \(K(G,1)\); the type \(F_\infty\) property alone does not establish either finiteness condition requested above. Thus the supplied assessment must not be treated as a verified resolution of the displayed questions.
\subsection{Short English statement}
This record combines two separate questions: whether a nonamenable group without nonabelian free subgroups can have a finite classifying space, or at least finite virtual cohomological dimension; and whether certain relations for the Brin--Navas group transfer to Thompson's group \(F\). The first formulation depends on an unresolved interpretation of ``free subgroups,'' and the second lacks essential definitions. The supplied partially solved label and dated literature assessment remain unverified curation.
\subsection{Sources}
\begin{itemize}
\item Supplied local UnsolvedMath registry record, permanent ID \(30003850\), title ``Nonamenable Groups without Free Subgroups,'' source problem code \texttt{catalogue entry}, in the collection ``Oberwolfach Reports Open Problems.'' The supplied proposed year is 2018; no proposer is identified. The record labels the problem \texttt{partially\_solved}, but also marks its statement \texttt{unrecoverable} and supplies no clean statement.
\item Supplied report citation: \emph{Cohomological and Metric Properties of Groups of Homeomorphisms of \(\R\)} (2019), Oberwolfach Reports, pp.\ 1579--1633, DOI \texttt{10.4171/owr/2018/26}. \url{https://doi.org/10.4171/owr/2018/26}. The background describes these reports as documenting mathematical workshops and their open-problem discussions.
\item Supplied statement-verification note: the official report is said to show that the extraction merges two independent numbered questions concerning nonamenable groups and Wajnryb--Witowicz relations. No separate clean statement was promoted, and the original record was retained for provenance. This source-verification claim is supplied metadata, not an independently checked finding here.
\item Supplied literature source: Lodha, \emph{A nonamenable type \(F_\infty\) group}. \url{https://arxiv.org/abs/1408.3127}.
\item Supplied literature source: Guba, \emph{How to solve equations in groups?} \url{https://gcc.episciences.org/11315/pdf}.
\item Supplied literature source: \emph{Oberwolfach Report 26/2018}. \url{https://ems.press/content/serial-article-files/46748?nt=1}.
\item Literature assessment dated exactly \texttt{2026-08-21}: the supplied curation attributes to the Lodha--Moore construction a nonamenable group of type \(F_\infty\) with no nonabelian free subgroups and describes this as settling the existence question; it also says that the Wajnryb--Witowicz question about \(F\) remains open after a proposed proof failed. This is supplied dated curation derived from the OWR material, not an independently verified literature assessment. Its claimed resolution does not, from the properties stated in this record alone, establish the stronger finiteness conditions in the extracted question. The source explicitly requests independent verification of time-sensitive status and cited claims.
\end{itemize}
\end{document}
